\documentclass[11pt,a4paper]{article}
% Common preamble for the hint sheets (\input from each file)
\usepackage[utf8]{inputenc}
\usepackage[T1]{fontenc}
\usepackage{lmodern}
\usepackage{amsmath,amssymb,amsthm,mathtools}
\usepackage[a4paper,margin=2.3cm]{geometry}
\usepackage{enumitem}
\usepackage{xcolor}
\usepackage[hidelinks]{hyperref}
\renewcommand{\leq}{\leqslant}
\renewcommand{\le}{\leqslant}
\renewcommand{\geq}{\geqslant}
\renewcommand{\ge}{\geqslant}
\setlength{\parindent}{0pt}
\setlength{\parskip}{0.35em}
% \hintstitle{Sheet title}{Class work / Homework}
\newcommand{\hintstitle}[2]{\begin{center}{\LARGE\bfseries Hints: #1\par}\vspace{0.5em}{\large #2\par}\end{center}
\noindent\textit{The hints to each problem are cumulative. Read one, try again, and open the next only if you are still stuck.}\par\medskip\hrule\medskip}
\newcommand{\problem}[1]{\par\bigskip{\large\bfseries Problem #1}\par\nopagebreak\smallskip}
\newcommand{\subpart}[1]{\par\smallskip\textbf{(#1)}\par\nopagebreak}
\newcommand{\hint}[1]{\par\smallskip\noindent\textcolor{black!70}{\textbf{Hint #1.}}\enspace}
\newcommand{\note}[1]{\par\smallskip\noindent\textit{Note.} #1\par}
\providecommand{\lt}{<}
\providecommand{\gt}{>}
\providecommand{\tightlist}{\setlength{\itemsep}{0pt}\setlength{\parskip}{0pt}}
\emergencystretch=3em

\begin{document}
\hintstitle{Abstract algebra -- 6. The hardest problems}{Homework}

\problem{1}
\hint{1} $G_p$ is an increasing union of the subgroups $\mu_{p^n}$, hence a subgroup. A subgroup $K\le G_p$ either contains elements of arbitrarily large $p$-power order (then $K\supseteq\mu_{p^n}$ for all $n$, so $K=G_p$) or lies in some $\mu_{p^N}$ (then it is finite).
\hint{2} Conversely, let $H$ be infinite with all proper subgroups finite. An element of infinite order would generate an infinite cyclic subgroup, which has the infinite proper subgroup $\langle x^2\rangle$. So $H$ consists of roots of unity.
\hint{3} Split $H$ into its $q$-parts $H_q$ ($q$ prime). If infinitely many $H_q$ are non-trivial, dropping one gives an infinite proper subgroup. So only finitely many are non-trivial, and some $H_p$ is infinite, so $H_p=H$. An infinite $p$-subgroup of the roots of unity contains $\mu_{p^n}$ for every $n$, so $H=G_p$.

\problem{2}
\note{``Ring'' means ring with unit, as everywhere on these sheets; with non-unital rings one would get $i(\mathbb Z_{12})=12$.}
\subpart{a}
\hint{1} A biadditive multiplication on $\mathbb Z_n$ is determined by $a=1\cdot1$: $x\cdot y=axy$. It is automatically associative and commutative. It has a unit iff $a$ is invertible mod $n$. So $i(\mathbb Z_n)=\varphi(n)$, and $\varphi(12)=4$.
\subpart{b}
\hint{1} Product rings: a pair of ring structures on $A$ and $B$ gives one on $A\times B$, and different pairs give different multiplications.
\subpart{c}
\hint{1} $H_k=\mathbb Z_{n_k}$ with $n_k=p_1p_2\cdots p_k$: $n_k/\varphi(n_k)=\prod_{i\le k}\frac{p_i}{p_i-1}\to\infty$, since $\sum1/p$ diverges.
\hint{2} $G_k=(\mathbb Z_2)^k$: transport the ring structure of $\mathbb F_2^k$ (or of $\mathbb F_{2^k}$) along all group automorphisms $\varphi\in\mathrm{GL}_k(\mathbb F_2)$. Two of them give the same multiplication only if they differ by a ring automorphism, and there are at most $k!$ (resp.\ $k$) of those. So $i(G_k)\ge|\mathrm{GL}_k(\mathbb F_2)|/k!$, which is about $2^{k^2}/k!$, far larger than $2^k$.

\problem{3}
\hint{1} Put $c=f(1)$ and suppose $c$ is even. Take $h=e$: you get $f(n)=f\bigl(n/\gcd(n,c)\bigr)$ for every $n$. In particular $f(2)=f(1)=c$. (If $c=0$, $f$ is already constant.)
\hint{2} Use a dihedral group: $g$ a rotation of order $a$, $h$ a reflection. Then $g^2h^{f(a)}$ is a reflection if $f(a)$ is odd and equals $g^2$ if $f(a)$ is even, while $h^ag^{f(2)}=h^ag^c$.
\hint{3} For even $a$: if $f(a)$ were odd you would get $c=f(2)=f(a/\gcd(a,c))=f(a)$, contradicting parity. So $f(a)$ is even, and then $f(a/2)=f(a)$. Hence $f(a)=f(2a)$ is even for every $a$.
\hint{4} For odd $a$ (with $f(a)$ even) the same group gives $f(a)=f(2)=c$. So $f$ is constant on odd numbers, and by halving on all numbers, a contradiction.

\problem{4}
\hint{1} The sum of all elements of a finite abelian group of exponent $p$ is $0$ unless the group is $\mathbb Z/2$. So with at least $4$ elements the non-zero elements add up to $0$. Then for $B$ you need $\sum_{x\in A}x=-\sum_{y\in B}y$.
\hint{2} Try $|B|=2$, $B=\{b,c\}$: you need $bc=-(b+c)$, i.e.\ $(1+b)(1+c)=1$.
\hint{3} Take $t\notin\{0,1,-1\}$ (possible since there are at least $4$ elements), $b=t-1$, $c=t^{-1}-1$. Then $b,c$ are non-zero and distinct, they commute, and $A=K^*\setminus B$ is non-empty.

\problem{5}
\subpart{a}
\hint{1} Linearise $[x,f(x)]=0$ by putting $x+y$.
\hint{2} For the second identity, compute $[x,f(xy)]=[x,f(x)f(y)]$ in two ways: directly, and via the first identity as $[f(x),xy]$. You get $(f(x)-x)[f(x),y]=0$. Rewrite $[f(x),y]=[x,f(y)]$ and use surjectivity of $f$.
\subpart{b}
\hint{1} In a division ring, (a) says: for each $x$, either $f(x)=x$ or $x$ commutes with everything. So $R$ is the union of two additive subgroups, $\{x:f(x)=x\}$ and $Z(R)$.
\hint{2} A group is never the union of two proper subgroups. Since $f\ne\mathrm{id}$, $Z(R)=R$.

\problem{6}
\hint{1} The identity reads $x(1-x)^2(1+x)=0$. Put $x=2$: $6=0$, and $2$ is invertible, so $3=0$.
\hint{2} If $n^2=0$, the identity gives $n=0$. So $A$ is reduced. In the commutative subring generated by $x$, $(x^3-x)^2$ is a multiple of $x(x-1)^2(x+1)$, so $x^3-x$ is nilpotent, hence $0$. Thus $x^3=x$ for all $x$.
\hint{3} Then $x^2$ is idempotent, and idempotents of a reduced ring are central. So $(x+1)^2-x^2-1=2x$ is central, and so is $x$: $A$ is commutative. Units satisfy $u^2=1$.
\hint{4} From $(a+b)^2=1=a^2$ get $b(2a+b)=0$. If $b$ were a unit, then $b=-2a=a$. And $(1+ab)^2=1+b(2a+b)=1$. Example: $\mathbb F_3\times\mathbb F_3$.

\problem{7}
\hint{1} Sum all the products: in an abelian group $\prod_i(a_ia_{i+1})=(\prod a_i)^2$ must equal $\prod a_i$. So the sum of all elements must vanish. For $\mathbb Z_n$ that sum is $n(n-1)/2$, which is $\ne0$ mod $n$ for even $n$.
\hint{2} For odd $n$ take $a_i=i$ ($i=0,\dots,n-1$): the sums $2i+1$ run over all residues, because $2$ is invertible.
\hint{3} For even order, look at a group with more than one involution, e.g.\ $\mathbb Z_2\times\mathbb Z_4$. The order $(0,0),(1,0),(0,1),(1,2),(1,1),(1,3),(0,3),(0,2)$ works.

\problem{8}
\hint{1} Polynomials modulo $x^p-x$ correspond to functions $\mathbb F_p\to\mathbb F_p$, and composition corresponds to composition. Identify the two generators as functions.
\hint{2} $x+1$ is the cycle $k\mapsto k+1$. For the other polynomial $q$: if $x\ne0,1$ then $q(x)=\frac{x^{p-1}-1}{x-1}+x=x$; also $q(0)=1$ and $q(1)=0$. So $q$ is the transposition $(0\,1)$.
\hint{3} A $p$-cycle and an adjacent transposition generate $S_p$. In a finite group the generated semigroup is the generated group. So $|W|=p!$.

\problem{9}
\hint{1} $n$ is interesting iff $x^n=1$ in $R=\mathbb F_p[x]/(x^p-x+1)$. So the smallest interesting number is the multiplicative order of $x$ in $R$.
\hint{2} In $R$, $x^p=x-1$, and by the Frobenius $x^{p^k}=x-k$. So $x^{p^p}=x$, and $x$ is invertible, hence $x^{p^p-1}=1$.
\hint{3} Norm trick: $x^{1+p+\dots+p^{p-1}}=\prod_{k\in\mathbb F_p}(x-k)=x^p-x=-1$. So $x^{2(p^p-1)/(p-1)}=1$, which is a smaller interesting number when $p>3$.
\hint{4} For $p=2$: $x$ has order $3$. For $p=3$: the order divides $26$, $x^2\ne1$ and $x^{13}=-1\ne1$, so it is $26$. Answer: $p\in\{2,3\}$.

\problem{10}
\subpart{a}
\hint{1} First show $xa,ax\in D$ for $a\in D$: if $ax$ were a unit, $a$ would have a right inverse, and then $a^2=0$ forces $a=0$.
\hint{2} Show $D$ is closed under addition. If $a+b=u$ is a unit with $a,b\in D$, then $c=u^{-1}a\in D$ and $1-c=u^{-1}b\in D$, so $0=(1-c)^2=1-2c$ and $c$ would be a unit.
\hint{3} Now $a+xa\in D$, so $0=(a+xa)^2=a^2+axa+xa^2+(xa)^2=axa$.
\subpart{b}
\hint{1} Start with any non-zero $a\in D$. If $ab\ne0$ for some $b\in D$, replace $a$ by $a'=ab$. Then $a'D\subsetneq aD$: $ab\in aD$, but $ab\notin abD$, since $ab=abd$ would give $ab=abd^2=0$. Also $|Da'|=|Dab|\le|Da|$.
\hint{2} Symmetrically for $ba\ne0$. So $|aD|+|Da|$ strictly decreases while $a$ stays non-zero. By finiteness the process stops at a non-zero $a$ with $aD=Da=0$.

\end{document}
